\begin{abstract}
This work presents a quantitative comparison of two serving pipelines for an
XGBoost model on tabular data: FastAPI with ONNX Runtime, and BentoML with
\textit{adaptive batching}. A 219-cell open-loop benchmark sweep over
(workers, threads, RPS) was executed with the Vegeta load generator, and the
end-to-end latency was decomposed into three measurement layers (pure
inference, host HTTP, containerised deployment) to attribute overhead. To
compare a load-shedding system against a queueing system fairly, the headline
metric is goodput within a 50~ms SLA. FastAPI sustained between 99.7\% and
100\% of goodput across the matrix from 50 to 1000 RPS, while BentoML
collapsed between 350 and 600 RPS scaling with worker count, with the P99
latency growing by two orders of magnitude over the same RPS range. Direct
CPU-utilisation evidence collected on a 60-run mechanism subset confirmed
that the single asynchronous dispatcher coroutine per BentoML worker is the
bottleneck: BentoML consumed approximately 2.5 times the CPU of FastAPI under
identical load. For cheap single-row tree-ensemble inference on CPU,
adaptive batching's per-request plumbing exceeds any batching amortisation.
\end{abstract}

\begin{resumo}
Este trabalho apresenta uma comparação quantitativa entre dois pipelines de
\textit{serving} para um modelo XGBoost aplicado a dados tabulares: FastAPI
com ONNX Runtime e BentoML com \textit{adaptive batching}. Foi executada uma
varredura de modelo aberto de 219 células sobre (\textit{workers},
\textit{threads}, RPS) com o gerador de carga Vegeta, decompondo a latência
ponta a ponta em três camadas de medição (inferência pura, HTTP no host e
\textit{deployment} conteinerizado) para atribuir o sobrecusto. Para comparar
de forma justa um sistema que descarta requisições contra um sistema que as
enfileira, a métrica principal adotada foi a vazão útil dentro de um SLA de
50~ms. O FastAPI sustentou entre 99,7\% e 100\% de vazão útil sobre toda a
matriz de 50 a 1000 RPS, enquanto o BentoML colapsou entre 350 e 600 RPS
escalando com o número de \textit{workers}, e a latência P99 cresceu duas
ordens de magnitude sobre o mesmo intervalo de RPS. A evidência direta de
utilização de CPU coletada em um subconjunto de mecanismo de 60 execuções
confirmou que a única corrotina assíncrona de \textit{dispatcher} por
\textit{worker} do BentoML é o gargalo: o BentoML consumiu cerca de 2,5 vezes
a CPU do FastAPI sob carga idêntica. Para inferência barata de árvores em
CPU sobre uma linha por requisição, o sobrecusto por requisição do
\textit{adaptive batching} excede qualquer amortização de lote.

\textbf{Palavras-chave:} serviços de inferência de \textit{ML}; latência de
cauda; P99; FastAPI; BentoML; ONNX Runtime; \textit{adaptive batching};
benchmarking de modelo aberto; \textit{Coordinated Omission}; XGBoost; dados
tabulares.
\end{resumo}

